\documentclass[12pt,a4paper,oneside]{report}

%================================================
% 1. PACKAGES
%================================================

\usepackage[utf8]{inputenc}
\usepackage[T1]{fontenc}
\usepackage[french]{babel}
\usepackage{lmodern}
\usepackage{geometry}
\usepackage{graphicx}
\usepackage{float}
\usepackage{array}
\usepackage{booktabs}
\usepackage{longtable}
\usepackage{tabularx}
\usepackage{multirow}
\usepackage{amsmath,amssymb}
\usepackage{xcolor}
\usepackage{setspace}
\usepackage{hyperref}
\usepackage{fancyhdr}
\usepackage{enumitem}
\usepackage{titlesec}
\usepackage{caption}
\usepackage{microtype}

%================================================
% 2. MISE EN PAGE
%================================================

\geometry{
    left=2.5cm,
    right=2.5cm,
    top=2.5cm,
    bottom=2.5cm
}

\onehalfspacing

\setlength{\parindent}{0.8cm}
\setlength{\parskip}{4pt}

\hypersetup{
    colorlinks=true,
    linkcolor=black,
    citecolor=blue,
    urlcolor=blue,
    pdftitle={Détection des maladies des feuilles de pomme de terre},
    pdfauthor={Rapport de projet}
}

\pagestyle{fancy}
\fancyhf{}
\fancyhead[L]{Détection des maladies des feuilles de pomme de terre}
\fancyhead[R]{Rapport de projet}
\fancyfoot[C]{\thepage}
\setlength{\headheight}{15pt}

\renewcommand{\chaptermark}[1]{
    \markboth{\thechapter.\ #1}{}
}

\titleformat{\chapter}[display]
{\normalfont\bfseries\Huge}
{\chaptertitlename\ \thechapter}
{15pt}
{\Huge}

\renewcommand{\arraystretch}{1.3}

%================================================
% 3. DEBUT DU DOCUMENT
%================================================

\begin{document}

%================================================
% PAGE DE GARDE
%================================================

\begin{titlepage}
\begin{center}

{\Large République Tunisienne}\\[0.3cm]

{\large Ministère de l'Enseignement Supérieur
et de la Recherche Scientifique}\\[1cm]

{\Large \textbf{Rapport de Projet}}\\[1cm]

\rule{\textwidth}{1pt}\\[0.4cm]

{\Huge \textbf{Détection des maladies des feuilles
de pomme de terre par Deep Learning}}\\[0.4cm]

{\Large Classification d'images à l'aide d'un
réseau de neurones convolutif (CNN)}\\[0.4cm]

\rule{\textwidth}{1pt}\\[1cm]

% Ajouter le logo de l'établissement si disponible
% \includegraphics[width=3cm]{logo.png}\\[1cm]

\textbf{Réalisé par :}\\[0.2cm]
Nom et prénom de l'étudiante 1\\
Nom et prénom de l'étudiante 2\\[0.6cm]

\textbf{Encadré par :}\\
Nom de l'encadrant\\[1cm]

\textbf{Année universitaire : 2026--2027}

\end{center}
\end{titlepage}

%================================================
% PAGES PRELIMINAIRES
%================================================

\pagenumbering{roman}

\chapter*{Remerciements}
\addcontentsline{toc}{chapter}{Remerciements}

Nous tenons à exprimer notre profonde gratitude à toutes
les personnes qui ont contribué, directement ou indirectement,
à la réalisation de ce projet.

Nous remercions particulièrement notre encadrant pour
son accompagnement, ses conseils et ses orientations
tout au long de ce travail.

Nous adressons également nos remerciements à l'ensemble
de nos enseignants pour les connaissances et les compétences
qu'ils nous ont transmises durant notre formation.

Enfin, nous remercions toutes les personnes qui nous ont
soutenus et encouragés dans la réalisation de ce projet.

%------------------------------------------------

\chapter*{Résumé}
\addcontentsline{toc}{chapter}{Résumé}

Les maladies des plantes constituent un problème important
pour le secteur agricole, car elles peuvent provoquer une
diminution de la qualité et du rendement des cultures.
La pomme de terre est particulièrement exposée à plusieurs
maladies foliaires, notamment l'Early Blight et le Late Blight.

L'identification de ces maladies repose souvent sur
l'observation visuelle des feuilles. Cependant, cette méthode
peut être difficile et nécessite des connaissances spécifiques
en agriculture.

Dans ce contexte, ce projet vise à développer une approche
automatique de classification des maladies des feuilles
de pomme de terre à partir d'images, en utilisant les techniques
du Deep Learning.

La solution étudiée repose sur un réseau de neurones
convolutif (CNN), développé avec TensorFlow et Keras.
Le modèle permet de distinguer trois catégories :
les feuilles atteintes par l'Early Blight, les feuilles atteintes
par le Late Blight et les feuilles saines.

Le processus comprend plusieurs étapes : la préparation
des images, leur redimensionnement, leur normalisation,
l'application de techniques de Data Augmentation,
l'entraînement du CNN et l'évaluation de ses performances.

Une étude de l'existant permet également de comparer
notre approche avec plusieurs travaux scientifiques récents
portant sur la classification des maladies des plantes.

\textbf{Mots-clés :} Intelligence artificielle, Deep Learning,
CNN, classification d'images, PlantVillage, pomme de terre,
Data Augmentation.

%------------------------------------------------

\chapter*{Abstract}
\addcontentsline{toc}{chapter}{Abstract}

Plant diseases represent a major challenge for agriculture
because they can reduce crop quality and productivity.
Potato plants are particularly affected by several leaf diseases,
including Early Blight and Late Blight.

This project focuses on the automatic classification of
potato leaf diseases using image processing and Deep Learning.

The proposed approach is based on a Convolutional Neural
Network (CNN) implemented with TensorFlow and Keras.
The model classifies leaf images into three categories:
Early Blight, Late Blight and Healthy.

The main stages include image preprocessing, resizing,
normalization, data augmentation, model training and
performance evaluation.

\textbf{Keywords:} Artificial Intelligence, Deep Learning,
CNN, Image Classification, PlantVillage, Potato Diseases.

%------------------------------------------------

\tableofcontents
\listoffigures
\listoftables

\clearpage
\pagenumbering{arabic}

%================================================
% INTRODUCTION GENERALE
%================================================

\chapter*{Introduction générale}
\addcontentsline{toc}{chapter}{Introduction générale}

L'agriculture constitue un secteur essentiel pour le
développement économique et la sécurité alimentaire.
Cependant, la production agricole est confrontée à plusieurs
difficultés, parmi lesquelles les maladies des plantes
occupent une place importante.

Les maladies végétales peuvent affecter différentes parties
des plantes, notamment les feuilles, les tiges et les racines.
Lorsqu'elles ne sont pas détectées à temps, elles peuvent
provoquer des pertes de rendement et une diminution de
la qualité des produits agricoles.

La pomme de terre (\textit{Solanum tuberosum}) représente
une culture importante dans de nombreux pays. Elle peut
être touchée par différentes maladies foliaires, en particulier
l'Early Blight et le Late Blight. Ces maladies présentent
des symptômes visibles sur les feuilles, tels que des taches,
des changements de couleur et des zones nécrosées.

Traditionnellement, le diagnostic repose sur l'observation
visuelle effectuée par les agriculteurs ou les spécialistes.
Toutefois, cette méthode dépend de l'expérience de la personne
et peut devenir difficile lorsque les symptômes sont similaires.

Avec les progrès de l'intelligence artificielle, et plus
particulièrement du Deep Learning, de nouvelles solutions
permettent d'automatiser l'analyse des images.
Les réseaux de neurones convolutifs (CNN) sont notamment
utilisés pour reconnaître les caractéristiques visuelles
et classifier automatiquement les images.

Dans ce contexte, notre projet porte sur la détection
et la classification des maladies des feuilles de pomme
de terre à partir d'images en utilisant un modèle CNN.

Le travail s'appuie sur un projet de référence disponible
sur GitHub, consacré à la classification des maladies
de la pomme de terre avec TensorFlow et Keras.

Notre démarche consiste à étudier et adapter cette approche,
à comprendre les différentes étapes de son fonctionnement,
à préparer les données et à analyser les performances
du modèle.

Le rapport est organisé en six chapitres.
Le premier présente la problématique et les objectifs.
Le deuxième est consacré à l'étude de l'existant et aux
travaux scientifiques récents.
Le troisième présente les méthodologies CRISP-DM et
CRISP-ML(Q). Le quatrième décrit le dataset et sa préparation.
Le cinquième présente la conception et l'entraînement du CNN.
Enfin, le sixième chapitre est consacré à l'évaluation
des performances et à l'analyse des résultats.

%================================================
% CHAPITRE 1
%================================================

\chapter{Problématique et objectifs}

\section{Contexte du projet}

L'agriculture joue un rôle fondamental dans la satisfaction
des besoins alimentaires et dans le développement économique.
Toutefois, les cultures agricoles sont exposées à plusieurs
maladies susceptibles d'affecter leur croissance et leur
productivité.

La pomme de terre est particulièrement concernée par
différentes maladies foliaires. Parmi les plus fréquentes,
on distingue l'Early Blight et le Late Blight.

L'Early Blight se manifeste généralement par l'apparition
de taches brunes sur les feuilles, souvent caractérisées
par des anneaux concentriques. Le Late Blight peut provoquer
des zones sombres et des lésions qui se développent rapidement,
notamment dans des conditions environnementales favorables
au développement de la maladie.

La détection précoce de ces maladies est importante afin
de limiter leur propagation et de réduire les pertes agricoles.

Cependant, l'identification manuelle des symptômes peut
être difficile, particulièrement lorsque les maladies
présentent des caractéristiques visuelles proches.

L'intelligence artificielle offre une possibilité intéressante
pour automatiser cette tâche. Grâce à la vision par ordinateur
et aux réseaux de neurones convolutifs, il devient possible
d'analyser les images des feuilles et d'apprendre les
caractéristiques associées à chaque maladie.

Notre projet s'inscrit donc dans le domaine de l'agriculture
intelligente et vise à utiliser le Deep Learning pour
faciliter la reconnaissance des maladies des feuilles
de pomme de terre.

\section{Problématique}

Le diagnostic des maladies des plantes constitue une tâche
importante, mais qui présente plusieurs difficultés.

Premièrement, l'identification manuelle nécessite une
connaissance des symptômes et de leurs différences.
Deuxièmement, certaines maladies peuvent présenter des
signes visuels similaires, ce qui augmente le risque
d'erreur de diagnostic.

De plus, l'analyse d'un grand nombre d'images peut demander
beaucoup de temps. Une solution automatique pourrait donc
faciliter cette opération et aider à obtenir un premier
diagnostic à partir d'une photographie.

La problématique principale de notre projet est formulée
comme suit :

\begin{quote}
\textit{Comment concevoir un système intelligent basé sur
le Deep Learning, capable d'identifier automatiquement
les maladies des feuilles de pomme de terre à partir
d'images et de distinguer les feuilles saines des feuilles
atteintes par différentes maladies ?}
\end{quote}

\section{Objectif général}

L'objectif général de ce projet est de développer et d'étudier
un système de classification automatique des maladies
des feuilles de pomme de terre en utilisant un réseau
de neurones convolutif (CNN).

Le système doit être capable d'analyser une image en entrée
et de déterminer la catégorie correspondante parmi trois
classes : Early Blight, Late Blight et Healthy.

\section{Objectifs spécifiques}

Pour atteindre cet objectif général, plusieurs objectifs
spécifiques ont été définis :

\begin{itemize}
    \item Étudier les maladies foliaires de la pomme de terre
    et leurs caractéristiques visuelles.
    \item Explorer un dataset d'images contenant des feuilles
    saines et malades.
    \item Organiser les images selon leurs catégories.
    \item Effectuer le prétraitement des images.
    \item Appliquer des techniques de Data Augmentation
    afin d'augmenter la diversité des données d'entraînement.
    \item Construire un modèle de classification basé sur
    un réseau de neurones convolutif.
    \item Entraîner le modèle à reconnaître les trois classes.
    \item Évaluer les performances à l'aide de métriques
    adaptées à la classification.
    \item Étudier les possibilités d'utilisation du modèle
    dans une application de diagnostic agricole.
\end{itemize}

\section{Résultats attendus}

À la fin du projet, nous cherchons à obtenir un modèle
capable de classifier automatiquement les images des
feuilles de pomme de terre.

Les résultats attendus sont les suivants :

\begin{itemize}
    \item Un dataset organisé et préparé pour l'apprentissage.
    \item Un modèle CNN entraîné sur les images disponibles.
    \item Une classification en trois catégories.
    \item Une évaluation quantitative du modèle.
    \item Une analyse des erreurs de classification.
    \item Une base pour le développement d'un système
    d'aide au diagnostic des maladies des plantes.
\end{itemize}

Il est important de préciser que les performances réelles
du modèle seront déterminées à partir des expériences
d'entraînement et de test. Elles ne peuvent pas être
annoncées avant l'exécution et l'évaluation du modèle.

%================================================
% CHAPITRE 2
%================================================

\chapter{État de l'art et bibliographie}

\section{Introduction}

La détection automatique des maladies des plantes à partir
d'images représente un domaine de recherche important
en intelligence artificielle appliquée à l'agriculture.

Les premières approches reposaient principalement sur
le traitement d'images et l'extraction manuelle de
caractéristiques visuelles. Avec le développement du
Machine Learning et du Deep Learning, les chercheurs
ont progressivement proposé des méthodes capables
d'apprendre automatiquement les caractéristiques des images.

Dans le cas des maladies de la pomme de terre, plusieurs
travaux récents se sont intéressés à la reconnaissance
de l'Early Blight, du Late Blight et des feuilles saines.

Cette partie présente trois études scientifiques récentes
qui permettent de comprendre les principales approches
existantes et de positionner notre travail par rapport
à ces recherches.

\section{Première étude : Talab et al. (2024)}

\subsection{Présentation de l'étude}

Dans leur article intitulé :

\textit{Investigating the accuracy of classification in
unbalanced data in order to diagnose two common potato
leaf diseases (early blight and late blight) using image
processing and machine learning},

Talab, Mohammadzamani et Parashkoohi (2024) étudient
la classification des maladies des feuilles de pomme
de terre en présence d'un déséquilibre entre les classes.

Les auteurs s'intéressent à trois catégories :
Early Blight, Late Blight et Healthy.

\subsection{Méthode utilisée}

Contrairement aux approches basées directement sur
l'apprentissage profond, cette étude utilise une méthode
fondée sur le traitement d'images et l'extraction
de caractéristiques.

Les auteurs commencent par isoler les régions affectées
des feuilles, puis extraient plusieurs types de caractéristiques :

\begin{itemize}
    \item 45 caractéristiques de couleur.
    \item 99 caractéristiques de texture.
    \item 6 caractéristiques de forme.
\end{itemize}

Les caractéristiques sont extraites à partir de plusieurs
espaces colorimétriques, notamment RGB, Lab et HSV.

Ensuite, l'algorithme Random Forest est utilisé pour
classifier les images à partir des caractéristiques extraites.

Les auteurs étudient également l'effet du déséquilibre
des données en utilisant des techniques d'Over-Sampling
et d'Under-Sampling.

\subsection{Résultats obtenus}

Les résultats rapportés dans l'article montrent que
l'équilibrage des données peut améliorer les performances
de classification.

Les précisions obtenues sont présentées dans le tableau
suivant :

\begin{table}[H]
\centering
\caption{Résultats de Talab et al. (2024)}
\begin{tabular}{lc}
\toprule
\textbf{Configuration} & \textbf{Accuracy (\%)}\\
\midrule
Dataset original & 87.89\\
Over-Sampling & 91.09\\
Under-Sampling & 96.17\\
\bottomrule
\end{tabular}
\end{table}

\subsection{Limites de l'étude}

Cette approche nécessite une étape d'extraction manuelle
des caractéristiques avant la classification.

La qualité des caractéristiques sélectionnées peut donc
influencer directement les performances du système.

De plus, l'utilisation d'un algorithme classique de
Machine Learning dépend de la capacité des caractéristiques
extraites à représenter correctement les différences
entre les maladies.

\section{Deuxième étude : Alhammad et al. (2025)}

\subsection{Présentation de l'étude}

Dans leur article intitulé :

\textit{Deep learning and explainable AI for classification
of potato leaf diseases},

Alhammad et ses collaborateurs proposent une approche
basée sur le Deep Learning pour classifier les maladies
des feuilles de pomme de terre.

Le travail combine l'utilisation d'un modèle pré-entraîné
et des techniques d'Intelligence Artificielle Explicable
(Explainable Artificial Intelligence - XAI).

L'objectif est non seulement d'obtenir une classification,
mais également de mieux comprendre les régions de l'image
qui influencent la décision du modèle.

\subsection{Méthode utilisée}

Les auteurs utilisent un dataset dérivé de PlantVillage,
composé de 2\,152 images réparties en trois catégories :

\begin{itemize}
    \item 1\,000 images Early Blight.
    \item 1\,000 images Late Blight.
    \item 152 images Healthy.
\end{itemize}

Afin de réduire les effets du déséquilibre entre les classes,
des techniques d'augmentation sont utilisées, notamment
la rotation, le changement d'échelle, le retournement
et la translation verticale.

La méthode repose sur le Transfer Learning, qui permet
d'utiliser les connaissances acquises par un réseau
pré-entraîné pour réaliser une nouvelle tâche de classification.

Les auteurs utilisent également Grad-CAM afin de visualiser
les régions de l'image ayant contribué à la prédiction.

\subsection{Résultats obtenus}

L'étude rapporte les résultats suivants :

\begin{itemize}
    \item Accuracy de validation : 97\%.
    \item Accuracy de test : 98\%.
\end{itemize}

L'utilisation de Grad-CAM permet également de fournir
une représentation visuelle des zones importantes
pour la classification.

\subsection{Limites de l'étude}

L'utilisation d'un modèle pré-entraîné peut améliorer
les performances, mais nécessite de gérer correctement
le choix de l'architecture et les paramètres de
fine-tuning.

Par ailleurs, l'utilisation de Grad-CAM améliore
l'interprétabilité, mais ne garantit pas à elle seule
que la prédiction soit toujours correcte.

\section{Troisième étude : Bhavani et Chalapathi (2026)}

\subsection{Présentation de l'étude}

Dans l'article intitulé :

\textit{PotatoLeafNet: two-stage convolutional neural
networks for effective Potato Leaf disease identification
and classification},

Bhavani et Chalapathi proposent une architecture de
Deep Learning dédiée à l'identification des maladies
des feuilles de pomme de terre.

Le travail porte sur les trois catégories principales :
Healthy, Early Blight et Late Blight.

\subsection{Méthode utilisée}

Les auteurs proposent une approche qui associe une
stratégie d'augmentation des images à un réseau CNN
compact composé de onze couches.

Le dataset étudié contient 4\,072 images réparties
en trois catégories.

Les images sont redimensionnées à 224 $\times$ 224 pixels
et normalisées dans l'intervalle [0,1].

La stratégie de Data Augmentation comprend plusieurs
transformations :

\begin{itemize}
    \item Rotation.
    \item Translation.
    \item Shear.
    \item Zoom.
    \item Retournement horizontal.
    \item Modification de la luminosité.
    \item Modification des canaux de couleur.
\end{itemize}

Ces transformations sont appliquées aux données
d'entraînement afin d'augmenter leur diversité.

Le modèle est développé avec TensorFlow et Keras,
et utilise l'optimiseur Adam ainsi qu'une fonction
de perte Categorical Crossentropy.

\subsection{Résultats obtenus}

Les auteurs rapportent les performances suivantes
sur le jeu de test :

\begin{itemize}
    \item Accuracy : 98.52\%.
    \item Macro-Precision : 98.67\%.
    \item Macro-Recall : 99.67\%.
    \item Macro-F1-score : 99.16\%.
\end{itemize}

Ces résultats montrent l'intérêt de combiner
une architecture CNN adaptée avec une stratégie
d'augmentation des données.

\subsection{Limites de l'étude}

Malgré les résultats obtenus sur le dataset utilisé,
la généralisation à des images prises directement
dans les champs reste un défi.

Les conditions réelles peuvent introduire des variations
de luminosité, des arrière-plans complexes, des feuilles
partiellement cachées et des différences entre les variétés
de pommes de terre.

Une validation complémentaire sur des images provenant
de conditions réelles est donc importante.

\section{Comparaison des études scientifiques}

Les trois études précédentes utilisent des approches
différentes pour résoudre le même problème général.

Le tableau suivant présente une comparaison synthétique.

\begin{table}[H]
\centering
\caption{Comparaison des travaux scientifiques}
\small
\begin{tabularx}{\textwidth}{p{2.5cm}p{2.3cm}X p{2cm}}
\toprule
\textbf{Étude} &
\textbf{Approche} &
\textbf{Caractéristiques principales} &
\textbf{Résultat rapporté}\\
\midrule

Talab et al. (2024) &
Random Forest &
Extraction de caractéristiques de couleur, texture
et forme. Étude du déséquilibre des classes. &
96.17\% avec Under-Sampling\\

Alhammad et al. (2025) &
Transfer Learning et XAI &
Utilisation d'un réseau pré-entraîné, augmentation
des images et visualisation Grad-CAM. &
98\% sur le test\\

Bhavani et Chalapathi (2026) &
CNN personnalisé &
CNN de onze couches et augmentation séquentielle
des images. &
98.52\% sur le test\\

\bottomrule
\end{tabularx}
\end{table}

\textit{Remarque : les performances présentées proviennent
des articles respectifs. Les datasets, les partitions
et les protocoles d'évaluation ne sont pas nécessairement
identiques. Les valeurs ne constituent donc pas une
comparaison expérimentale directe entre les modèles.}

\section{Analyse comparative}

L'analyse des travaux scientifiques permet de dégager
plusieurs observations.

Premièrement, les méthodes classiques de Machine Learning,
comme Random Forest, peuvent obtenir de bons résultats
lorsque les caractéristiques de couleur, de texture
et de forme sont correctement extraites.

Cependant, elles nécessitent une étape préalable de
conception et de sélection des caractéristiques.

Deuxièmement, les approches basées sur le Transfer Learning
permettent de réutiliser des caractéristiques apprises
sur de grandes bases d'images. Elles peuvent être
particulièrement utiles lorsque les données disponibles
sont limitées.

Troisièmement, les CNN personnalisés peuvent apprendre
directement les caractéristiques visuelles à partir
des images. L'ajout de techniques de Data Augmentation
peut améliorer la diversité des exemples présentés
pendant l'entraînement.

Les travaux étudiés montrent également que la qualité
et la diversité du dataset jouent un rôle important.
Un modèle performant sur des images prises dans un
environnement contrôlé ne garantit pas nécessairement
les mêmes performances dans des conditions agricoles réelles.

\section{Solutions commerciales existantes}

En complément des recherches scientifiques, plusieurs
applications commerciales utilisent l'intelligence
artificielle pour aider les agriculteurs à identifier
les problèmes affectant leurs cultures.

\subsection{Plantix}

Plantix est une application destinée au diagnostic
des problèmes rencontrés par les plantes.

Elle permet à l'utilisateur de prendre ou de télécharger
une photographie d'une plante afin d'obtenir une
identification du problème.

L'application propose également des informations sur
les symptômes, la prévention et les possibilités
de traitement.

Plantix illustre ainsi l'intérêt de combiner la
reconnaissance d'images avec des informations agricoles
destinées à accompagner l'utilisateur.

\subsection{Agrio}

Agrio est une solution numérique destinée à la surveillance
et à la protection des cultures.

Elle propose notamment l'identification des maladies,
des ravageurs et des problèmes nutritionnels à partir
de photographies.

La plateforme fournit également des suggestions
d'intervention et des outils de suivi des cultures.

\subsection{Comparaison des solutions commerciales}

\begin{table}[H]
\centering
\caption{Comparaison des solutions commerciales}
\begin{tabularx}{\textwidth}{p{2.5cm}X X}
\toprule
\textbf{Critère} & \textbf{Plantix} & \textbf{Agrio}\\
\midrule
Diagnostic par image &
Oui &
Oui\\

Identification des maladies &
Oui &
Oui\\

Informations et conseils &
Symptômes, prévention et traitements &
Identification et suggestions de solutions\\

Suivi agricole &
Informations et communauté &
Suivi des cultures, alertes et outils agricoles\\

Objectif principal &
Aider à identifier et traiter les problèmes des plantes &
Accompagner la surveillance et la protection des cultures\\
\bottomrule
\end{tabularx}
\end{table}

Ces solutions montrent que la valeur d'un système agricole
intelligent ne se limite pas à reconnaître une maladie.
Il est également important de présenter des informations
compréhensibles et des recommandations utiles.

\section{Positionnement de notre projet}

Après l'étude des travaux scientifiques et des solutions
commerciales, nous avons choisi de nous orienter vers
une approche basée sur un réseau de neurones convolutif
(CNN).

Notre projet s'appuie sur une architecture CNN développée
avec TensorFlow et Keras. Le modèle apprend directement
à partir des images et permet de classifier les feuilles
en trois catégories.

Contrairement à l'approche de Talab et al. (2024), qui
nécessite l'extraction de caractéristiques manuelles
avant l'utilisation de Random Forest, notre approche
utilise un réseau capable d'apprendre automatiquement
les caractéristiques visuelles.

Par rapport à l'étude d'Alhammad et al. (2025), notre
travail adopte une architecture CNN personnalisée,
sans présenter le Transfer Learning ou Grad-CAM comme
des éléments déjà intégrés dans notre implémentation.

Enfin, les travaux de Bhavani et Chalapathi (2026)
montrent l'intérêt des CNN personnalisés et de
l'augmentation des données, ce qui rejoint les principes
de notre démarche.

\subsection{Valeur ajoutée de notre projet}

La valeur ajoutée de notre projet repose sur plusieurs
éléments complémentaires :

\begin{enumerate}
    \item \textbf{Une approche simple et reproductible :}
    l'utilisation d'un CNN construit avec TensorFlow
    et Keras permet de disposer d'une architecture
    compréhensible et réutilisable.

    \item \textbf{Une classification ciblée :}
    le système se concentre sur trois catégories importantes
    pour la pomme de terre : Early Blight, Late Blight
    et Healthy.

    \item \textbf{Une préparation des images adaptée :}
    le redimensionnement, la normalisation et la
    Data Augmentation permettent de préparer les images
    et d'augmenter la diversité des données d'entraînement.

    \item \textbf{Une interprétation plus utile de la prédiction :}
    au-delà du nom de la classe, une évolution possible
    consiste à afficher le niveau de confiance associé
    à la prédiction et à signaler les cas incertains.

    \item \textbf{Une aide à la décision agricole :}
    une extension envisagée consiste à associer chaque
    maladie détectée à une fiche contenant les symptômes,
    les mesures préventives et des recommandations
    agricoles documentées.
\end{enumerate}

Il convient de distinguer les fonctionnalités réellement
implémentées des améliorations proposées.

Le modèle CNN et la préparation des images constituent
la base technique du travail. Le module de recommandations
et la gestion des prédictions incertaines représentent
une extension fonctionnelle à développer et à valider.

Notre contribution ne consiste donc pas à proposer
une nouvelle architecture de réseau de neurones, mais
à construire une solution de classification ciblée,
à étudier son fonctionnement et à proposer son
enrichissement par des informations utiles à l'utilisateur.

%================================================
% CHAPITRE 3
%================================================

\chapter{Méthodologie CRISP-DM et CRISP-ML(Q)}

\section{Introduction}

La réalisation d'un projet de Machine Learning nécessite
une démarche organisée permettant de passer de la
compréhension du problème à la construction et à
l'évaluation du modèle.

Dans ce projet, nous nous appuyons principalement sur
la méthodologie CRISP-DM et sur les principes de
CRISP-ML(Q).

Ces méthodologies permettent de structurer les différentes
étapes du développement et de prendre en compte
la qualité des données, la qualité du modèle et son
suivi après son déploiement.

\section{Présentation de CRISP-DM}

CRISP-DM signifie \textit{Cross-Industry Standard Process
for Data Mining}.

Il s'agit d'une méthodologie organisée en six phases
permettant de structurer un projet de traitement
et d'analyse des données.

Les six phases sont :

\begin{enumerate}
    \item Business Understanding.
    \item Data Understanding.
    \item Data Preparation.
    \item Modeling.
    \item Evaluation.
    \item Deployment.
\end{enumerate}

\section{Présentation de CRISP-ML(Q)}

CRISP-ML(Q) est une méthodologie destinée aux projets
de Machine Learning qui met l'accent sur la qualité
tout au long du cycle de vie du modèle.

Elle prend en compte les différentes étapes de conception,
de préparation, d'entraînement, d'évaluation, de déploiement
et de maintenance.

Son intérêt est de ne pas considérer l'entraînement
comme la seule étape importante. La qualité des données,
la reproductibilité, le suivi et la surveillance du modèle
doivent également être pris en considération.

\section{Justification du choix des méthodologies}

Le choix de CRISP-DM est justifié par sa structure
progressive et sa capacité à organiser les différentes
étapes d'un projet basé sur les données.

Dans notre cas, il permet de passer de la compréhension
du problème agricole à l'exploration du dataset, puis
à la préparation des images et à l'entraînement du CNN.

CRISP-ML(Q) complète cette démarche en mettant l'accent
sur la qualité du processus et la surveillance du modèle.

L'utilisation de ces deux méthodologies permet donc
d'obtenir une démarche structurée et adaptée à notre
projet de classification d'images.

\section{Application de CRISP-DM à notre projet}

\subsection{Compréhension du besoin}

La première phase consiste à comprendre le problème
à résoudre.

Notre objectif est de faciliter l'identification des
maladies des feuilles de pomme de terre à partir d'images.

Le système doit distinguer trois catégories :
Early Blight, Late Blight et Healthy.

\subsection{Compréhension des données}

Cette phase consiste à étudier les données disponibles.

Nous utilisons un ensemble d'images consacré aux maladies
des feuilles de pomme de terre, issu du dataset PlantVillage.

L'exploration permet de comprendre :

\begin{itemize}
    \item La structure des dossiers.
    \item Le nombre d'images disponibles.
    \item Les différentes catégories.
    \item La distribution des images entre les classes.
    \item Les caractéristiques visuelles des feuilles.
\end{itemize}

\subsection{Préparation des données}

Cette étape consiste à préparer les images pour leur
utilisation dans le réseau CNN.

Elle comprend le redimensionnement, la normalisation
et l'application de transformations d'augmentation
aux données d'entraînement.

\subsection{Modélisation}

La phase de modélisation consiste à construire le réseau
de neurones convolutif et à définir ses paramètres.

Le modèle est développé avec TensorFlow et Keras.
Il apprend à reconnaître les caractéristiques visuelles
associées aux trois classes.

\subsection{Évaluation}

L'évaluation permet de mesurer la qualité des prédictions
du modèle à l'aide de plusieurs métriques, notamment
l'Accuracy, la Precision, le Recall et le F1-score.

Une matrice de confusion permet également d'analyser
les erreurs de classification entre les catégories.

\subsection{Déploiement}

La dernière phase concerne l'utilisation du modèle
dans un environnement permettant de réaliser des
prédictions sur de nouvelles images.

Le dépôt GitHub de référence présente une architecture
plus complète qui comprend notamment une API, une
interface Web et une application mobile.

Dans le cadre de notre travail, les composants effectivement
réalisés et testés doivent être distingués des composants
présents uniquement dans le dépôt de référence.

\section{Application des principes CRISP-ML(Q)}

Les principes de CRISP-ML(Q) peuvent être appliqués
à notre projet de la manière suivante :

\begin{itemize}
    \item \textbf{Qualité des données :} vérifier les
    étiquettes, les dimensions et la cohérence des images.

    \item \textbf{Qualité de la préparation :} s'assurer
    que les transformations ne modifient pas la classe
    réelle de la feuille.

    \item \textbf{Qualité du modèle :} surveiller
    l'apprentissage et détecter un éventuel surapprentissage.

    \item \textbf{Qualité de l'évaluation :} utiliser
    des données indépendantes pour mesurer les performances.

    \item \textbf{Qualité du déploiement :} vérifier
    que les nouvelles images sont correctement préparées
    avant la prédiction.

    \item \textbf{Maintenance :} prévoir la mise à jour
    du dataset et le réentraînement du modèle lorsque
    de nouvelles données représentatives sont disponibles.
\end{itemize}

\section{Cycle de vie et maintenance du modèle}

Un modèle de Deep Learning ne doit pas être considéré
comme une solution définitive.

Les conditions agricoles peuvent évoluer et de nouvelles
variétés de plantes ou de nouveaux symptômes peuvent
apparaître.

Il est donc nécessaire de prévoir :

\begin{itemize}
    \item L'ajout de nouvelles images.
    \item La vérification des étiquettes.
    \item Le suivi des performances.
    \item L'analyse des erreurs de prédiction.
    \item Le réentraînement lorsque cela est nécessaire.
\end{itemize}

Ces opérations contribuent à maintenir la qualité
et la pertinence du système.

%================================================
% CHAPITRE 4
%================================================

\chapter{Présentation et préparation du dataset}

\section{Présentation du dataset PlantVillage}

Le dataset utilisé dans notre projet est consacré
à la reconnaissance des maladies des plantes à partir
d'images de feuilles.

Le projet de référence indique le téléchargement des
données depuis Kaggle, puis la conservation des dossiers
correspondant à la pomme de terre.

La version classique du sous-ensemble Potato de PlantVillage
utilisée dans plusieurs études comprend 2\,152 images
réparties en trois classes.

La distribution est la suivante :

\begin{table}[H]
\centering
\caption{Répartition des images du dataset Potato}
\begin{tabular}{lc}
\toprule
\textbf{Classe} & \textbf{Nombre d'images}\\
\midrule
Potato Early Blight & 1\,000\\
Potato Late Blight & 1\,000\\
Potato Healthy & 152\\
\midrule
\textbf{Total} & \textbf{2\,152}\\
\bottomrule
\end{tabular}
\end{table}

\textit{Important : cette répartition correspond à la
version classique de 2\,152 images. Le nombre exact doit
être confirmé par le comptage des fichiers présents
dans le dataset réellement utilisé pour l'entraînement.}

\section{Structure du dataset}

Les images sont organisées en dossiers correspondant
aux différentes classes.

Chaque dossier contient les images associées à une
catégorie précise.

Cette organisation facilite le chargement automatique
des images et l'attribution des étiquettes nécessaires
à l'apprentissage supervisé.

\section{Présentation des classes}

\subsection{Potato Early Blight}

L'Early Blight est une maladie foliaire qui peut provoquer
des taches brunes sur les feuilles.

Ces taches peuvent présenter des anneaux concentriques
caractéristiques. La reconnaissance de ces motifs
visuels constitue un élément important pour la classification.

\subsection{Potato Late Blight}

Le Late Blight est une maladie qui peut entraîner
l'apparition de lésions sombres sur les feuilles.
Son développement peut être rapide lorsque les conditions
environnementales sont favorables.

\subsection{Potato Healthy}

Cette catégorie regroupe les images de feuilles saines,
ne présentant pas les symptômes caractéristiques
des deux maladies précédentes.

\section{Analyse exploratoire des images}

L'analyse exploratoire permet de visualiser les images
et de comprendre les différences entre les catégories.

Elle permet notamment d'observer :

\begin{itemize}
    \item Les différences de couleur.
    \item Les formes des taches.
    \item Les textures présentes sur les feuilles.
    \item Les variations d'orientation.
    \item Les différences de taille et de qualité des images.
\end{itemize}

% Ajouter ici des exemples d'images des trois classes.
%
% \begin{figure}[H]
% \centering
% \includegraphics[width=0.8\textwidth]{classes.png}
% \caption{Exemples des trois catégories du dataset}
% \end{figure}

\section{Répartition des données}

Pour entraîner et évaluer le réseau, les données doivent
être organisées en ensembles distincts.

\begin{itemize}
    \item \textbf{Training :} utilisé pour apprendre
    les paramètres du réseau.
    \item \textbf{Validation :} utilisé pour suivre
    l'apprentissage et contrôler la généralisation.
    \item \textbf{Test :} utilisé pour mesurer les
    performances finales sur des images non utilisées
    pendant l'entraînement.
\end{itemize}

La répartition exacte doit suivre les paramètres
du notebook utilisé.

Il est important de conserver une séparation indépendante
des données afin d'éviter que des images utilisées
pendant l'entraînement se retrouvent dans le jeu de test.

\section{Prétraitement des images}

Avant leur utilisation dans le CNN, les images doivent
être transformées afin de respecter un format commun.

\subsection{Redimensionnement}

Les images sont redimensionnées à une taille fixe
de 256 $\times$ 256 pixels.

Cette opération permet de garantir que toutes les images
présentent les mêmes dimensions et peuvent être regroupées
dans des lots pendant l'entraînement.

\subsection{Normalisation}

Les valeurs des pixels sont initialement comprises
entre 0 et 255.

La normalisation consiste à diviser chaque valeur
par 255 afin d'obtenir des valeurs comprises entre
0 et 1.

La formule utilisée est :

\[
X_{\text{normalisé}} = \frac{X}{255}
\]

Cette transformation facilite le traitement numérique
des images et contribue à stabiliser l'apprentissage.

\section{Data Augmentation}

\subsection{Définition}

La Data Augmentation est une technique qui permet
d'augmenter la diversité des images d'entraînement
en appliquant différentes transformations aux images
existantes.

L'objectif est de présenter au réseau des variations
des mêmes exemples afin qu'il apprenne à reconnaître
les caractéristiques importantes indépendamment
de certaines transformations visuelles.

Cette technique peut contribuer à limiter le surapprentissage
et à améliorer la capacité de généralisation du modèle.

\subsection{RandomFlip}

RandomFlip est une transformation qui permet de retourner
aléatoirement les images horizontalement et verticalement.

Elle aide le réseau à apprendre que l'orientation d'une
feuille ne doit pas nécessairement modifier sa catégorie.

\subsection{RandomRotation}

RandomRotation applique une rotation aléatoire aux images.

Dans le code de référence, le facteur de rotation
est fixé à 0.2.

Cette transformation permet de présenter au réseau
des feuilles ayant différentes orientations.

\subsection{Implémentation de la Data Augmentation}

Le code utilisé dans le notebook repose sur les couches
de prétraitement de TensorFlow/Keras.

\begin{verbatim}
data_augmentation = tf.keras.Sequential([
    layers.experimental.preprocessing.RandomFlip(
        "horizontal_and_vertical"
    ),
    layers.experimental.preprocessing.RandomRotation(0.2),
])
\end{verbatim}

Les transformations sont appliquées aux images
d'entraînement.

Les données de validation et de test doivent rester
indépendantes des transformations aléatoires utilisées
pour enrichir l'entraînement.

\section{Synthèse de la préparation des données}

Les différentes opérations réalisées avant l'entraînement
sont résumées dans le tableau suivant.

\begin{table}[H]
\centering
\caption{Synthèse des opérations de préparation}
\begin{tabularx}{\textwidth}{p{3cm}X}
\toprule
\textbf{Opération} & \textbf{Description}\\
\midrule
Chargement &
Lecture des images organisées par classe.\\

Redimensionnement &
Conversion des images à une taille de 256 $\times$ 256 pixels.\\

Normalisation &
Transformation des valeurs des pixels dans l'intervalle [0,1].\\

RandomFlip &
Retournement horizontal et vertical aléatoire.\\

RandomRotation &
Rotation aléatoire des images.\\

Organisation &
Préparation des ensembles d'entraînement, de validation
et de test.\\
\bottomrule
\end{tabularx}
\end{table}

%================================================
% CHAPITRE 5
%================================================

\chapter{Modélisation et implémentation du CNN}

\section{Introduction}

Après la préparation du dataset, l'étape suivante consiste
à construire un modèle capable d'apprendre les caractéristiques
visuelles des feuilles de pomme de terre.

Pour réaliser cette tâche, nous utilisons un réseau
de neurones convolutif (CNN).

Ce type de réseau est particulièrement adapté au traitement
des images, car il permet d'extraire automatiquement
des caractéristiques telles que les contours, les textures,
les formes et les motifs visuels.

Le modèle est développé avec TensorFlow et Keras.

\section{Présentation du modèle CNN}

Un réseau de neurones convolutif est une architecture
de Deep Learning conçue pour traiter des données
structurées sous forme de grilles, notamment les images.

Contrairement aux méthodes qui nécessitent une extraction
manuelle des caractéristiques, le CNN apprend automatiquement
les représentations nécessaires à partir des images
d'entraînement.

Le réseau est organisé en plusieurs couches qui réalisent
successivement l'extraction des caractéristiques et
la classification.

\section{Architecture générale du CNN}

Le modèle utilisé dans le notebook repose sur une
architecture séquentielle de type Sequential.

Il comprend des couches de convolution et de pooling,
suivies de couches entièrement connectées pour réaliser
la classification.

Le modèle reçoit une image de taille 256 $\times$ 256
avec trois canaux de couleur RGB.

L'architecture comprend notamment :

\begin{itemize}
    \item Des couches Conv2D pour extraire les caractéristiques.
    \item Des couches MaxPooling2D pour réduire les dimensions.
    \item Une couche Flatten pour transformer les
    caractéristiques en vecteur.
    \item Une couche Dense pour réaliser la classification.
    \item Une couche de sortie contenant trois neurones.
\end{itemize}

\section{Description des couches du réseau}

\subsection{Couche de convolution Conv2D}

La couche Conv2D applique des filtres sur l'image
afin de détecter des caractéristiques visuelles.

Les premières couches peuvent apprendre à détecter
des contours, des couleurs et des textures simples.
Les couches suivantes permettent d'apprendre des motifs
plus complexes.

La convolution peut être représentée par :

\[
Y(i,j)=\sum_m\sum_n X(i-m,j-n)K(m,n)
\]

où $X$ représente l'image d'entrée, $K$ le filtre
de convolution et $Y$ la carte de caractéristiques obtenue.

\subsection{Couche MaxPooling2D}

La couche MaxPooling2D réduit les dimensions spatiales
des cartes de caractéristiques.

Elle conserve les valeurs maximales dans des régions
locales afin de réduire le nombre de calculs.

Cette opération permet également de rendre le réseau
moins sensible à certaines petites variations de position.

\subsection{Couche Flatten}

La couche Flatten transforme les cartes de caractéristiques
en un vecteur unidimensionnel.

Cette transformation permet de transmettre les
informations extraites vers les couches entièrement
connectées.

\subsection{Couche Dense}

La couche Dense est une couche entièrement connectée.

Elle combine les caractéristiques extraites afin
de déterminer les informations nécessaires à la
classification.

Le modèle de référence comprend une couche Dense
de 64 neurones avec une fonction d'activation ReLU.

La fonction ReLU est définie par :

\[
f(x)=\max(0,x)
\]

\subsection{Fonction Softmax}

La couche de sortie utilise la fonction Softmax afin
de produire une probabilité pour chaque classe.

La formule est :

\[
P(y=i)=\frac{e^{z_i}}{\sum_{j=1}^{K}e^{z_j}}
\]

Dans notre projet, $K=3$, car le système distingue
trois catégories.

La classe ayant la probabilité la plus élevée
correspond à la prédiction du modèle.

\section{Résumé de l'architecture}

Le tableau suivant présente les principaux composants
de l'architecture étudiée.

\begin{table}[H]
\centering
\caption{Principales couches du modèle CNN}
\begin{tabularx}{\textwidth}{p{3.5cm}X}
\toprule
\textbf{Couche} & \textbf{Rôle}\\
\midrule
Resizing &
Adapter les dimensions des images.\\

Rescaling &
Normaliser les valeurs des pixels.\\

Conv2D &
Extraire les caractéristiques visuelles.\\

MaxPooling2D &
Réduire les dimensions des cartes de caractéristiques.\\

Flatten &
Transformer les caractéristiques en vecteur.\\

Dense &
Combiner les caractéristiques pour la classification.\\

Softmax &
Produire les probabilités associées aux trois classes.\\
\bottomrule
\end{tabularx}
\end{table}

\section{Configuration du modèle}

Le modèle est construit avec TensorFlow et Keras.

Les principaux paramètres de configuration sont présentés
dans le tableau suivant.

\begin{table}[H]
\centering
\caption{Configuration générale du modèle}
\begin{tabular}{ll}
\toprule
\textbf{Paramètre} & \textbf{Valeur}\\
\midrule
Type de modèle & CNN\\
Framework & TensorFlow / Keras\\
Taille des images & 256 $\times$ 256\\
Nombre de canaux & 3 (RGB)\\
Batch size & 32\\
Nombre de classes & 3\\
Optimiseur & Adam\\
Fonction de perte & Sparse Categorical Crossentropy\\
Nombre d'époques & 50 (configuration de référence)\\
\bottomrule
\end{tabular}
\end{table}

\textit{Les valeurs doivent être vérifiées dans la version
finale du notebook exécuté. Si les paramètres ont été
modifiés pendant les expérimentations, le tableau doit
être mis à jour.}

\section{Entraînement du CNN}

L'entraînement consiste à présenter progressivement
les images au réseau afin qu'il apprenne à associer
chaque image à sa catégorie.

Au début, les paramètres du réseau sont initialisés.
Le modèle réalise ensuite des prédictions et compare
ses sorties avec les étiquettes réelles.

L'erreur calculée est utilisée pour ajuster les poids
du réseau à l'aide de l'algorithme de rétropropagation.

Ce processus est répété pendant plusieurs époques
jusqu'à ce que le modèle apprenne à reconnaître
les caractéristiques associées aux différentes classes.

\section{Fonction de perte}

La fonction de perte mesure la différence entre les
prédictions du modèle et les catégories réelles.

Dans le projet, la fonction Sparse Categorical
Crossentropy est utilisée pour la classification
multi-classes lorsque les étiquettes sont représentées
par des entiers.

Elle peut être exprimée par :

\[
L=-\frac{1}{N}\sum_{i=1}^{N}
\log(P(y_i))
\]

où $P(y_i)$ représente la probabilité attribuée
par le modèle à la classe correcte de l'exemple $i$.

\section{Optimiseur Adam}

Adam est un algorithme d'optimisation utilisé pour
mettre à jour les paramètres du réseau pendant
l'entraînement.

Il adapte les taux d'apprentissage à partir des
estimations des moyennes et des variances des gradients.

Son utilisation permet de faciliter la convergence
du réseau et d'améliorer le processus d'apprentissage.

\section{Sauvegarde du modèle}

Après l'entraînement, le modèle peut être sauvegardé
afin de pouvoir être réutilisé sans devoir recommencer
tout le processus d'apprentissage.

Le dépôt GitHub de référence contient notamment
des modèles sauvegardés au format H5 ainsi que
des éléments destinés à la conversion TensorFlow Lite.

\section{Prédiction sur une nouvelle image}

Une fois le modèle entraîné, une nouvelle image peut
être fournie au système.

Elle doit subir les mêmes opérations de prétraitement
que les images d'entraînement, notamment le redimensionnement
et la normalisation.

Le modèle calcule ensuite les probabilités associées
aux trois catégories.

La classe correspondant à la probabilité la plus élevée
est retournée comme résultat de classification.

\[
\hat{y}=\arg\max_{i\in\{1,2,3\}}P(y=i)
\]

Le résultat peut être présenté sous la forme suivante :

\begin{itemize}
    \item Potato Early Blight.
    \item Potato Late Blight.
    \item Potato Healthy.
\end{itemize}

%================================================
% CHAPITRE 6
%================================================

\chapter{Évaluation des performances}

\section{Objectif de l'évaluation}

L'évaluation représente une étape essentielle dans
un projet de classification.

Elle permet de mesurer la capacité du modèle à reconnaître
correctement les différentes catégories et à généraliser
sur des images qui n'ont pas été utilisées pendant
l'entraînement.

Dans notre projet, l'évaluation doit permettre de vérifier
si le CNN distingue correctement les feuilles saines
des feuilles atteintes par l'Early Blight et le Late Blight.

\section{Métriques d'évaluation}

Plusieurs métriques sont utilisées pour évaluer
les performances d'un modèle de classification.

Dans un problème multi-classes, ces métriques peuvent
être calculées pour chaque classe, puis agrégées.

\subsection{Accuracy}

L'Accuracy représente la proportion de prédictions
correctes parmi l'ensemble des prédictions.

\[
Accuracy=\frac{\text{Nombre de prédictions correctes}}
{\text{Nombre total de prédictions}}
\]

Une Accuracy élevée indique que le modèle réalise
un grand nombre de prédictions correctes.

\subsection{Precision}

La Precision mesure la proportion des prédictions
positives qui sont réellement positives.

\[
Precision=\frac{TP}{TP+FP}
\]

Elle permet d'évaluer la capacité du modèle à éviter
les fausses alertes pour une classe donnée.

\subsection{Recall}

Le Recall mesure la proportion des exemples réellement
positifs qui ont été correctement identifiés.

\[
Recall=\frac{TP}{TP+FN}
\]

Cette métrique est particulièrement importante dans
la détection des maladies, car un faux négatif peut
correspondre à une feuille malade classée comme saine.

\subsection{F1-score}

Le F1-score représente la moyenne harmonique entre
Precision et Recall.

\[
F1=2\times
\frac{Precision\times Recall}
{Precision+Recall}
\]

Il permet de prendre en compte simultanément
la précision et la capacité de détection.

\subsection{Matrice de confusion}

La matrice de confusion permet de comparer les
classes réelles avec les classes prédites par le modèle.

Elle permet d'identifier les erreurs entre les différentes
catégories.

Pour notre projet, elle contient trois lignes et
trois colonnes correspondant aux classes suivantes :

\begin{itemize}
    \item Early Blight.
    \item Late Blight.
    \item Healthy.
\end{itemize}

\begin{table}[H]
\centering
\caption{Structure de la matrice de confusion}
\begin{tabular}{lccc}
\toprule
\multirow{2}{*}{\textbf{Classe réelle}} &
\multicolumn{3}{c}{\textbf{Classe prédite}}\\
\cmidrule(lr){2-4}
& Early Blight & Late Blight & Healthy\\
\midrule
Early Blight & & &\\
Late Blight & & &\\
Healthy & & &\\
\bottomrule
\end{tabular}
\end{table}

Les cases de la diagonale représentent les prédictions
correctes. Les autres cases représentent les erreurs
de classification.

\section{Visualisation de l'apprentissage}

L'évolution de l'Accuracy et de la Loss pendant
l'entraînement permet d'observer le comportement
du réseau.

La courbe d'Accuracy permet de suivre l'évolution
de la proportion de prédictions correctes.

La courbe de Loss permet de suivre l'évolution
de l'erreur du modèle.

La comparaison entre les courbes d'entraînement
et de validation permet notamment de détecter
un éventuel surapprentissage.

% Ajouter les graphiques obtenus après l'entraînement.
%
% \begin{figure}[H]
% \centering
% \includegraphics[width=0.8\textwidth]{accuracy.png}
% \caption{Évolution de l'Accuracy}
% \end{figure}
%
% \begin{figure}[H]
% \centering
% \includegraphics[width=0.8\textwidth]{loss.png}
% \caption{Évolution de la fonction de perte}
% \end{figure}

\section{Analyse des résultats}

Les résultats doivent être interprétés à partir des
mesures réellement obtenues pendant les expérimentations.

L'analyse ne doit pas se limiter à l'Accuracy globale.
Il est également important d'étudier les performances
pour chacune des trois catégories.

Une attention particulière doit être accordée à la
classe Healthy, car elle contient moins d'images que
les deux catégories de maladies dans la version classique
du dataset de 2\,152 images.

Les résultats obtenus devront être présentés dans
le tableau suivant.

\begin{table}[H]
\centering
\caption{Résultats expérimentaux du modèle CNN}
\begin{tabular}{lc}
\toprule
\textbf{Métrique} & \textbf{Valeur obtenue}\\
\midrule
Accuracy & À compléter\\
Precision macro & À compléter\\
Recall macro & À compléter\\
F1-score macro & À compléter\\
\bottomrule
\end{tabular}
\end{table}

Aucune valeur ne doit être inventée. Les cases seront
complétées après l'exécution et l'évaluation du modèle
sur le jeu de test.

\section{Limites du modèle}

Malgré l'intérêt des réseaux CNN pour la classification
des images, plusieurs limites doivent être prises
en considération.

\begin{itemize}
    \item \textbf{Déséquilibre des classes :}
    la différence du nombre d'images entre les catégories
    peut influencer l'apprentissage.

    \item \textbf{Variations des conditions de prise de vue :}
    les images utilisées dans les datasets peuvent être
    différentes des photographies prises directement
    dans les champs.

    \item \textbf{Arrière-plans :}
    les feuilles photographiées dans des environnements
    complexes peuvent présenter des difficultés
    supplémentaires pour le modèle.

    \item \textbf{Similarité des symptômes :}
    certaines maladies peuvent présenter des caractéristiques
    visuelles proches.

    \item \textbf{Généralisation :}
    un bon résultat sur un dataset ne garantit pas
    nécessairement les mêmes performances sur des
    images provenant d'autres environnements.
\end{itemize}

\section{Perspectives d'amélioration}

Plusieurs améliorations peuvent être envisagées pour
rendre le système plus complet.

\begin{itemize}
    \item Enrichir le dataset avec des photographies
    prises dans des conditions agricoles réelles.

    \item Étudier l'effet de différentes stratégies
    de Data Augmentation.

    \item Ajouter une méthode d'explicabilité comme
    Grad-CAM afin de visualiser les régions importantes
    pour la prédiction.

    \item Ajouter un mécanisme de gestion des prédictions
    incertaines.

    \item Développer un module de recommandations
    agricoles basé sur des sources scientifiques fiables.

    \item Tester le modèle dans une application
    permettant à l'utilisateur de télécharger une image
    et d'obtenir une prédiction.
\end{itemize}

%================================================
% CONCLUSION GENERALE
%================================================

\chapter*{Conclusion générale}
\addcontentsline{toc}{chapter}{Conclusion générale}

Ce projet s'inscrit dans le domaine de l'intelligence
artificielle appliquée à l'agriculture et porte sur
la classification automatique des maladies des feuilles
de pomme de terre à partir d'images.

L'objectif principal était d'étudier et de mettre en place
une approche basée sur le Deep Learning afin de distinguer
les feuilles atteintes par l'Early Blight, le Late Blight
et les feuilles saines.

Dans un premier temps, nous avons étudié le contexte
du problème et les difficultés liées au diagnostic
manuel des maladies végétales.

Ensuite, l'étude de l'existant nous a permis d'examiner
plusieurs travaux scientifiques récents utilisant
des approches différentes, notamment Random Forest,
le Transfer Learning et les architectures CNN.

Cette comparaison a mis en évidence l'intérêt des
réseaux de neurones convolutifs pour l'analyse automatique
des caractéristiques visuelles des feuilles.

Nous avons également présenté la méthodologie CRISP-DM
et les principes de CRISP-ML(Q), qui permettent
d'organiser les étapes de développement d'un projet
de Machine Learning.

La partie expérimentale repose sur la préparation
des images, le redimensionnement, la normalisation
et l'utilisation de techniques de Data Augmentation,
notamment RandomFlip et RandomRotation.

Le modèle étudié est un réseau CNN développé avec
TensorFlow et Keras. Il a pour objectif de classifier
les images en trois catégories.

L'évaluation des performances permet de mesurer
la capacité du réseau à réaliser cette classification
et d'identifier ses éventuelles limites.

Enfin, une perspective importante consiste à enrichir
la classification par des informations complémentaires,
telles que le niveau de confiance de la prédiction
et des recommandations agricoles documentées.

Ce travail constitue ainsi une base pour le développement
de solutions intelligentes destinées à faciliter
l'identification des maladies des plantes et à soutenir
la prise de décision dans le domaine agricole.

%================================================
% BIBLIOGRAPHIE
%================================================

\begin{thebibliography}{99}

\bibitem{talab2024}
H. K. Talab, D. Mohammadzamani et M. G. Parashkoohi,
\textit{Investigating the accuracy of classification in
unbalanced data in order to diagnose two common potato
leaf diseases (early blight and late blight) using image
processing and machine learning},
Discover Applied Sciences, vol. 6, article 286, 2024.

DOI :
\url{https://doi.org/10.1007/s42452-024-05959-2}

\bibitem{alhammad2025}
S. M. Alhammad, D. S. Khafaga, W. M. El-hady,
F. M. Samy et K. M. Hosny,
\textit{Deep learning and explainable AI for classification
of potato leaf diseases},
Frontiers in Artificial Intelligence, 2025.

DOI :
\url{https://doi.org/10.3389/frai.2024.1449329}

\bibitem{bhavani2026}
G. D. Bhavani et M. M. V. Chalapathi,
\textit{PotatoLeafNet: two-stage convolutional neural
networks for effective Potato Leaf disease identification
and classification},
Frontiers in Artificial Intelligence, 2026.

DOI :
\url{https://doi.org/10.3389/frai.2025.1668839}

\bibitem{codebasics}
Codebasics,
\textit{Potato Disease Classification: Training, REST APIs
and Frontend to test},
GitHub Repository.

\url{https://github.com/codebasics/potato-disease-classification}

\bibitem{plantix}
Plantix,
\textit{Plantix - Crop Diagnosis and Treatment Information}.

\url{https://www.plantix.net/en/}

\bibitem{agrio}
Agrio,
\textit{AI Agronomist and Plant Disease Identification}.

\url{https://agrio.app/}

\bibitem{tensorflow}
TensorFlow,
\textit{TensorFlow Documentation}.

\url{https://www.tensorflow.org/}

\bibitem{keras}
Keras,
\textit{Keras Documentation}.

\url{https://keras.io/}

\end{thebibliography}

%================================================
% WEBOGRAPHIE
%================================================

\chapter*{Webographie}
\addcontentsline{toc}{chapter}{Webographie}

\begin{itemize}
    \item GitHub : \url{https://github.com/}
    \item Kaggle : \url{https://www.kaggle.com/}
    \item TensorFlow : \url{https://www.tensorflow.org/}
    \item Keras : \url{https://keras.io/}
    \item Plantix : \url{https://www.plantix.net/en/}
    \item Agrio : \url{https://agrio.app/}
\end{itemize}

\end{document}